\chapter{Hito 12: Despliegue del Dashboard en Streamlit Community Cloud}

\section{Objetivo}

El objetivo de este hito es publicar el dashboard macrofinanciero en Internet utilizando Streamlit Community Cloud.

Al finalizar este hito el estudiante será capaz de:

\begin{itemize}
\item Conectar un repositorio GitHub con Streamlit Community Cloud.
\item Configurar el despliegue automático.
\item Resolver errores comunes de dependencias.
\item Gestionar la base DuckDB en un entorno remoto.
\item Publicar una aplicación accesible mediante URL pública.
\end{itemize}

\section{Resultado esperado}

Al finalizar este hito el proyecto deberá estar disponible en una dirección similar a:

\begin{verbatim}
https://dwh-macrofinanciero-rd.streamlit.app
\end{verbatim}

o cualquier URL generada por Streamlit Community Cloud.

\section{Arquitectura final}

La arquitectura completa del proyecto será:

\begin{verbatim}
GitHub
|
V
Streamlit Cloud
|
V
Dashboard Web

DuckDB
|
V
Capa Mart
|
V
Visualizaciones
\end{verbatim}

\section{Paso 1: Verificar que el proyecto funciona localmente}

Antes de desplegar, ejecutar:

\begin{lstlisting}[language=bash]
python etl/run_pipeline.py
\end{lstlisting}

y luego:

\begin{lstlisting}[language=bash]
streamlit run app/streamlit_app.py
\end{lstlisting}

Verificar:

\begin{itemize}
\item El dashboard abre correctamente.
\item Los gráficos se visualizan.
\item Los KPIs aparecen.
\item No existen errores de conexión.
\end{itemize}

\section{Paso 2: Verificar estructura del repositorio}

La estructura mínima requerida será:

\begin{verbatim}
dwh_macrofinanciero_rd/

+-- app/
|   -- streamlit_app.py

+-- data/
|   -- warehouse/
|       -- macrofinanzas.duckdb

+-- etl/

+-- sql/

+-- requirements.txt

+-- README.md
\end{verbatim}

\section{Paso 3: Verificar requirements.txt}

Abrir:

\begin{verbatim}
requirements.txt
\end{verbatim}

Confirmar que incluye:

\begin{verbatim}
duckdb
pandas
streamlit
plotly
python-dotenv
\end{verbatim}

Si existen dudas:

\begin{lstlisting}[language=bash]
pip freeze > requirements.txt
\end{lstlisting}

\section{Paso 4: Decidir estrategia para DuckDB}

Existen dos alternativas.

\subsection{Opción A: Publicar la base DuckDB}

Subir al repositorio:

\begin{verbatim}
data/warehouse/macrofinanzas.duckdb
\end{verbatim}

Ventajas:

\begin{itemize}
\item Configuración sencilla.
\item La aplicación funciona inmediatamente.
\item Adecuada para proyectos educativos.
\end{itemize}

Desventajas:

\begin{itemize}
\item La base crece con el tiempo.
\item No es adecuada para producción.
\end{itemize}

\subsection{Opción B: Reconstruir la base automáticamente}

La aplicación ejecuta el pipeline cuando detecta que la base no existe.

Ventajas:

\begin{itemize}
\item Más profesional.
\item Reproducible.
\item No requiere almacenar la base en Git.
\end{itemize}

Desventajas:

\begin{itemize}
\item Mayor complejidad.
\item Tiempo de inicio más largo.
\end{itemize}

\section{Recomendación}

Para este proyecto se recomienda inicialmente:

\begin{quote}
Subir una versión pequeña de \texttt{macrofinanzas.duckdb}.
\end{quote}

Posteriormente se puede migrar a una estrategia de reconstrucción automática.

\section{Paso 5: Crear cuenta en Streamlit Community Cloud}

Entrar a:

\begin{verbatim}
https://share.streamlit.io
\end{verbatim}

Autenticarse utilizando:

\begin{itemize}
\item Cuenta GitHub.
\end{itemize}

\section{Paso 6: Crear nueva aplicación}

Seleccionar:

\begin{verbatim}
New App
\end{verbatim}

Completar:

\begin{itemize}
\item Repository
\item Branch
\item Main file path
\end{itemize}

\section{Paso 7: Configurar archivo principal}

Seleccionar:

\begin{verbatim}
app/streamlit_app.py
\end{verbatim}

como archivo principal.

\section{Paso 8: Configurar variables de entorno}

Si el proyecto utiliza:

\begin{verbatim}
.env
\end{verbatim}

agregar variables desde:

\begin{verbatim}
Settings
Secrets
\end{verbatim}

Ejemplo:

\begin{verbatim}
DUCKDB_PATH=data/warehouse/macrofinanzas.duckdb
\end{verbatim}

\section{Paso 9: Iniciar despliegue}

Presionar:

\begin{verbatim}
Deploy
\end{verbatim}

Streamlit realizará:

\begin{enumerate}
\item Clonado del repositorio.
\item Instalación de dependencias.
\item Ejecución de la aplicación.
\end{enumerate}

\section{Paso 10: Revisar logs}

Si la aplicación falla, revisar:

\begin{verbatim}
Manage App
Logs
\end{verbatim}

Errores frecuentes:

\begin{itemize}
\item Dependencias faltantes.
\item Ruta incorrecta de DuckDB.
\item Archivo principal incorrecto.
\item Error en consultas SQL.
\end{itemize}

\section{Paso 11: Resolver problema de rutas}

En local utilizamos:

\begin{lstlisting}[language=Python]
DB_PATH =
"data/warehouse/macrofinanzas.duckdb"
\end{lstlisting}

Para hacer la aplicación más robusta:

\begin{lstlisting}[language=Python]
from pathlib import Path

BASE_DIR = Path(**file**).resolve().parent.parent

DB_PATH = (
BASE_DIR
/
"data"
/
"warehouse"
/
"macrofinanzas.duckdb"
)
\end{lstlisting}

Esta técnica evita errores cuando la aplicación se ejecuta en servidores remotos.

\section{Paso 12: Crear página de error amigable}

Agregar:

\begin{lstlisting}[language=Python]
try:

```
resumen = run_query(
    """
    SELECT *
    FROM mart.vw_resumen_macrofinanciero
    """
)
```

except Exception as e:

```
st.error(
    f"Error al conectar con DuckDB: {e}"
)

st.stop()
```

\end{lstlisting}

Esto facilita el diagnóstico de errores en producción.

\section{Paso 13: Verificar funcionamiento remoto}

Comprobar:

\begin{itemize}
\item La aplicación abre.
\item Los KPIs aparecen.
\item Los gráficos funcionan.
\item Los filtros responden.
\item No existen errores visibles.
\end{itemize}

\section{Paso 14: Crear versión inicial pública}

Una vez validado:

\begin{itemize}
\item Compartir URL.
\item Agregar URL al README.
\item Agregar URL al CV.
\item Agregar URL a LinkedIn.
\end{itemize}

\section{Paso 15: Actualizar README}

Agregar sección:

\begin{verbatim}

## Aplicación en línea

https://xxxxx.streamlit.app
\end{verbatim}

\section{Paso 16: Validar reproducibilidad}

Clonar el repositorio en otra carpeta.

Ejecutar:

\begin{lstlisting}[language=bash]
pip install -r requirements.txt

streamlit run app/streamlit_app.py
\end{lstlisting}

Si funciona correctamente, el proyecto es reproducible.

\section{Buenas prácticas}

Durante este hito seguiremos estas reglas:

\begin{enumerate}
\item No subir credenciales.
\item Mantener requirements actualizado.
\item Documentar claramente el despliegue.
\item Verificar el dashboard después de cada cambio.
\item Utilizar rutas relativas.
\end{enumerate}

\section{Checklist de validación}

\begin{itemize}
\item[$\Box$] Proyecto publicado en GitHub.
\item[$\Box$] Streamlit Cloud configurado.
\item[$\Box$] Aplicación desplegada.
\item[$\Box$] URL pública generada.
\item[$\Box$] Dashboard funcional.
\item[$\Box$] README actualizado.
\item[$\Box$] Proyecto reproducible.
\end{itemize}

\section{Entregable}

Al finalizar este hito el estudiante deberá disponer de:

\begin{itemize}
\item Dashboard público.
\item URL accesible desde Internet.
\item Repositorio GitHub profesional.
\item Pipeline reproducible.
\item Proyecto completo de Data Engineering y BI.
\end{itemize}

\section{Proyecto completado}

Al concluir este hito se habrá construido una solución completa:

\begin{verbatim}
CSV
|
V
Python
|
V
DuckDB
|
V
staging
|
V
core
|
V
mart
|
V
Streamlit
|
V
Internet
\end{verbatim}

El estudiante habrá desarrollado un proyecto integral de ingeniería de datos, modelado dimensional, ETL y visualización analítica aplicados al contexto macrofinanciero.
